\section{Numerical method}\label[appendix]{app:numerics}

All three-dimensional experiments use a Fourier pseudo-spectral discretization \citep{canuto1988spectral}
of \eqref{eq:hyper} on $\T^3=(\R/2\pi\mathbb Z)^3$ with an $N\times N\times N$ grid. Fields are stored as
Fourier coefficients $\hat u(k)$ and derivatives are computed spectrally. Incompressibility is
enforced by the Leray projection $P(k)=I-k\otimes k/|k|^2$ for $k\ne0$.

\paragraph{De-aliasing and real fields.}
Products are formed in physical space. The de-aliasing mask keeps the modes with
$|k_x|,|k_y|,|k_z|<\tfrac23\lfloor N/2\rfloor$, with strict inequality \citep{orszag1971aliasing}.
A non-strict cutoff would let two retained modes alias onto a retained mode when $N$ is divisible
by three. Runs without de-aliasing still remove the ambiguous Nyquist faces, since an even grid
identifies $+N/2$ with $-N/2$ and odd derivative symbols cannot be represented there. Restriction
to a coarser grid also drops those faces, which preserves Hermitian symmetry, real fields and
nested restriction. Spectra that are not those of real fields, and nonfinite states or forces, are
rejected rather than silently repaired. With the strict cutoff the semi-discrete nonlinear term
conserves energy; the test suite checks this cancellation.

\paragraph{Time integration.}
The linear part $L(k)=-(\nu|k|^2+\mu|k|^{2\alpha})$ is integrated exactly. The classical fourth-order
Runge--Kutta method is applied to the integrating-factor variable $e^{-Lt}\hat u$ (the Lawson
scheme), with $E=e^{L\Delta t}$ and $E_{1/2}=e^{L\Delta t/2}$:
\[
\begin{aligned}
a&=E_{1/2}\bigl(\hat u+\tfrac{\Delta t}2n_1\bigr),\quad
b=E_{1/2}\hat u+\tfrac{\Delta t}2n_2,\quad
c=E\hat u+\Delta t\,E_{1/2}n_3,\\
\hat u^{\,\mathrm{new}}&=E\hat u+\tfrac{\Delta t}6\bigl(En_1+2E_{1/2}n_2+2E_{1/2}n_3+n_4\bigr),
\end{aligned}
\]
where $n_1,\dots,n_4$ are the nonlinear and forcing terms at the stage states $\hat u,a,b,c$ and the
corresponding stage times. Exact forced-linear and quadratic scalar solutions confirm fourth-order
convergence in the tests. Runs end exactly at the requested time, taking a shorter final step when
necessary, and record the terminal state.

\paragraph{Forcing.}
When forcing is on, a seeded random field is projected to divergence-free form, truncated to
$|k|\le k_{\mathrm{force}}$, de-aliased and rescaled to the configured root-mean-square amplitude. The
forcing is deterministic given the seed and configuration.

\paragraph{Diagnostics.}
The local energy fractions are amplitude invariant and are computed after rescaling, so very small
nonzero shells are not lost to underflow; an empty shell is assigned zero. Ranks $m_{95}$ are
counts of singular values of the uncentered snapshot matrix needed to reach $95\%$ of its squared
Frobenius norm. A sampled work-to-energy ratio is reported as a ratio over one sampled window, not
as an extremum over admissible histories. Each producer records the configuration that applies to it (for evolution
runs the grid, time step, final time, viscosities, exponent and forcing parameters, with separate
coefficients for comparison runs; for certificates, the thresholds) in a printed
\texttt{FIGURE\_CONFIG} line and in its \texttt{.npz} file.

\paragraph{Scope.}
These simulations are diagnostics for the certificates and hinge quantities of the scaffold.
Low-resolution runs are not evidence about continuum blow-up or regularity.
